\subsection{Estructuras de control}

\subsubsection{Reglas de uso}

\begin{enumerate}
  \item \textbf{Toda estructura de control (\texttt{if}, \texttt{else}, \texttt{for}, \texttt{foreach}, \texttt{while}, \texttt{do-while}) deberá delimitar su cuerpo con llaves \texttt{\{\}}, incluso cuando contenga una única instrucción.} Por ejemplo, en \texttt{EnemyCombatController} cada rama de decisión deberá escribirse con llaves explícitas aunque solo invoque un método. Omitir las llaves facilita que una instrucción añadida posteriormente quede fuera del alcance previsto sin generar un error de compilación. La regla de estilo de código \texttt{IDE0011} de .NET, que rige el uso de llaves en C\#, establece \texttt{true} como valor predeterminado de la opción \texttt{csharp\_prefer\_braces}, es decir, preferir llaves incluso para una sola línea de código (Microsoft, 2023).

  \item \textbf{Se preferirán las cláusulas de guarda con retorno anticipado frente a estructuras \texttt{if-else} profundamente anidadas.} En \texttt{EnemyCombatController}, un estado \texttt{Dead} deberá descartarse al inicio del método mediante un \texttt{return} temprano en lugar de envolver el resto de la lógica dentro de un \texttt{if} adicional. Reducir el anidamiento disminuye el número de rutas de lectura que debe seguir mentalmente quien revisa el código. McConnell señala que limitar la profundidad de anidamiento y salir de una rutina tan pronto como se determina que no hay más trabajo que realizar mejora la comprensibilidad del código (McConnell, 2004).

  \item \textbf{Se utilizará \texttt{switch} en lugar de cadenas de \texttt{if-else if} cuando se evalúen tres o más ramas mutuamente excluyentes sobre la misma variable o expresión.} Por ejemplo, la transición entre \texttt{EnemyState.Idle}, \texttt{EnemyState.Chasing} y \texttt{EnemyState.Attacking} deberá resolverse con un \texttt{switch} sobre \texttt{currentState} en vez de comparaciones sucesivas del mismo valor. McConnell recomienda utilizar una estructura de selección múltiple en lugar de una cadena de condicionales encadenados cuando varias ramas se excluyen mutuamente sobre la misma variable (McConnell, 2004); la documentación de C\# describe, en el mismo sentido, que las instrucciones \texttt{if} y \texttt{switch} seleccionan la ruta de ejecución del programa a partir del valor de una expresión (Microsoft, 2026a).

  \item \textbf{Todo bloque \texttt{switch} deberá incluir una rama \texttt{default} explícita que gestione los valores no contemplados.} En \texttt{EnemyCombatController}, un valor de \texttt{EnemyState} no manejado deberá provocar una excepción controlada en la rama \texttt{default} en lugar de finalizar el método de forma silenciosa. McConnell indica que las estructuras de selección deben contemplar de forma explícita los casos no previstos para evitar que el programa continúe en un estado inconsistente (McConnell, 2004). La documentación de C\# confirma que, si ningún patrón coincide con la expresión evaluada y no existe una rama \texttt{default}, el \texttt{switch} no ejecuta ninguna acción para ese valor (Microsoft, 2026a).

  \item \textbf{Queda prohibido el uso de la instrucción \texttt{goto} para alterar el flujo de control.} Las salidas o repeticiones de un bucle deberán expresarse mediante \texttt{return}, \texttt{continue}, \texttt{break} o mediante la reestructuración de la condición, nunca mediante saltos incondicionales a una etiqueta. Dijkstra argumentó que la instrucción de salto incondicional dificulta establecer una correspondencia clara entre el texto de un programa y los procesos que este ejecuta en el tiempo, degradando la capacidad de razonar sobre su corrección (Dijkstra, 1968).

  \item \textbf{Las condiciones booleanas se expresarán de forma positiva y se evitará negar comparaciones cuando exista una forma equivalente más clara.} Deberá preferirse \texttt{if (currentState == EnemyState.Dead)} sobre construcciones equivalentes que nieguen la condición y desplacen la lógica principal a la rama negativa. Martin sostiene que las expresiones booleanas afirmativas se leen con mayor facilidad que sus equivalentes negadas y reducen el esfuerzo cognitivo necesario para seguir el flujo del programa (Martin, 2008).

  \item \textbf{Se preferirá \texttt{foreach} sobre un \texttt{for} indexado cuando el algoritmo no requiera el índice del elemento.} Recorrer una lista de objetivos con \texttt{foreach} expresa directamente la intención de procesar cada elemento de la colección sin exponer detalles de indexación innecesarios. La documentación de C\# describe \texttt{foreach} como la instrucción que ejecuta su cuerpo para cada elemento de una colección, sin que el desarrollador deba declarar ni actualizar manualmente un índice (Microsoft, 2026b).

  \item \textbf{Todo bucle deberá contar con una condición de terminación verificable y su variable de control no deberá modificarse dentro del propio cuerpo del bucle.} En \texttt{ApplyDamageToTargets}, la colección recorrida y la condición de salida deben permitir determinar, con solo leer la cabecera del bucle, cuántas veces se ejecutará su cuerpo. McConnell recomienda que el punto de entrada y la condición de terminación de un bucle sean fácilmente verificables y advierte contra la modificación de la variable de control dentro del cuerpo, ya que esta práctica dificulta predecir el número de iteraciones (McConnell, 2004).
\end{enumerate}

\subsubsection{Con estándar}

\begin{lstlisting}[style=csharpstyle]
public sealed class EnemyCombatController
{
    private readonly List<int> targetIds;

    public EnemyCombatController(List<int> targetIds)
    {
        this.targetIds = targetIds;
    }

    public void UpdateState(EnemyState currentState)
    {
        if (currentState == EnemyState.Dead)
        {
            return;
        }

        switch (currentState)
        {
            case EnemyState.Idle:
                Patrol();
                break;
            case EnemyState.Chasing:
                ChaseTarget();
                break;
            case EnemyState.Attacking:
                ApplyDamageToTargets();
                break;
            default:
                throw new ArgumentOutOfRangeException(nameof(currentState), currentState, "Estado de enemigo no soportado.");
        }
    }

    private void ApplyDamageToTargets()
    {
        foreach (int targetId in targetIds)
        {
            if (targetId > 0)
            {
                RegisterHit(targetId);
            }
        }
    }

    private void RegisterHit(int targetId)
    {
    }

    private void Patrol()
    {
    }

    private void ChaseTarget()
    {
    }
}

public enum EnemyState
{
    Idle,
    Chasing,
    Attacking,
    Dead
}
\end{lstlisting}

El método \texttt{UpdateState} descarta el estado \texttt{Dead} mediante una cláusula de guarda con retorno anticipado en lugar de anidar el resto de la lógica dentro de un condicional adicional, y resuelve las restantes ramas mutuamente excluyentes con un \texttt{switch} que incluye una rama \texttt{default} explícita. Todas las estructuras de control utilizan llaves aunque su cuerpo sea breve, la condición del \texttt{if} interno se expresa en forma positiva y el recorrido de \texttt{targetIds} se realiza con \texttt{foreach} sin necesidad de gestionar manualmente un índice, sin ninguna instrucción \texttt{goto} en el flujo del método.

\subsubsection{Sin estándar}

\begin{lstlisting}[style=csharpstyle]
public sealed class EnemyCombatController
{
    private readonly List<int> targetIds;

    public EnemyCombatController(List<int> targetIds)
    {
        this.targetIds = targetIds;
    }

    public void UpdateState(EnemyState currentState)
    {
        if (currentState != EnemyState.Dead)
        {
            if (currentState == EnemyState.Idle)
                Patrol();
            else if (currentState == EnemyState.Chasing)
                ChaseTarget();
            else if (currentState == EnemyState.Attacking)
                ApplyDamageToTargets();
        }
    }

    private void ApplyDamageToTargets()
    {
        for (int i = 0; i < targetIds.Count; i++)
        {
            if (targetIds[i] <= 0)
                goto SkipTarget;

            RegisterHit(targetIds[i]);

        SkipTarget:
            continue;
        }
    }

    private void RegisterHit(int targetId)
    {
    }

    private void Patrol()
    {
    }

    private void ChaseTarget()
    {
    }
}

public enum EnemyState
{
    Idle,
    Chasing,
    Attacking,
    Dead
}
\end{lstlisting}

El método \texttt{UpdateState} incumple el estándar porque envuelve toda la lógica dentro de un \texttt{if (currentState != EnemyState.Dead)} negado en lugar de utilizar una cláusula de guarda con retorno anticipado, y resuelve tres ramas mutuamente excluyentes con una cadena \texttt{if-else if} sin rama por defecto en lugar de un \texttt{switch}, dejando además cualquier estado no contemplado sin manejo explícito. Las ramas del condicional omiten las llaves pese a contener una instrucción, contradiciendo la preferencia por delimitar siempre el cuerpo de las estructuras de control establecida en la regla \texttt{IDE0011} (Microsoft, 2023). En \texttt{ApplyDamageToTargets}, el uso de \texttt{goto SkipTarget} para omitir un objetivo reintroduce el salto incondicional que Dijkstra identificó como perjudicial para razonar sobre el flujo de un programa, en lugar de emplear directamente \texttt{continue} (Dijkstra, 1968). Estas decisiones contradicen las recomendaciones de preferir cláusulas de guarda, condiciones afirmativas y estructuras de selección explícitas (McConnell, 2004; Martin, 2008).
